\documentclass[12pt]{article}
\usepackage[top=1in, bottom=1in, left=.75in, right=.75in]{geometry}
\usepackage{amsmath}
\usepackage{fancyhdr}
\usepackage{graphicx}
\usepackage{txfonts}
\usepackage{multicol,coordsys}
\usepackage[scaled=0.86]{helvet}
\renewcommand{\emph}[1]{\textsf{\textbf{#1}}}
\usepackage{anyfontsize}
\usepackage[shortlabels]{enumitem}
% \usepackage{times}
% \usepackage[lf]{MinionPro}
\usepackage{tikz,pgfplots,mathrsfs}
%\def\degC{{}^\circ{\rm C}}
\def\ra{\rightarrow}
\usetikzlibrary{calc, backgrounds}
\pgfplotsset{compat = newest}
\newcommand{\blank}[1]{\rule{#1}{0.75pt}}

\pgfplotsset{my style/.append style={axis x line=middle, axis y line=
middle, xlabel={$x$}, ylabel={$y$},axis equal}}

%yticklabels={,,} , xticklabels={,,}

% \setmainfont{Times}
% \def\sansfont{Lucida Grande Bold}
\parindent 0pt
\parskip 4pt
\pagestyle{fancy}
\fancyfoot[C]{\emph{\thepage}}
%% version indicator
\fancyfoot[R]{\emph{v1}}
\fancyhead[L]{\ifnum \value{page} > 1\relax\emph{Math 252: Midterm Exam 2}\fi}
\fancyhead[R]{\ifnum \value{page} > 1\relax\emph{Fall 2025}\fi}
\headheight 15pt
\renewcommand{\headrulewidth}{0pt}
\renewcommand{\footrulewidth}{0pt}
\let\ds\displaystyle
\def\continued{{\emph {Continued....}}}
\def\continuing{{\emph {Problem \arabic{probcount} continued....}}\par\vskip 4pt}


\newcounter{probcount}
\newcounter{subprobcount}
\newcommand{\thesubproblem}{\emph{\alph{subprobcount}.}}
\def\problem#1{\setcounter{subprobcount}{0}%
\addtocounter{probcount}{1}{\emph{\arabic{probcount}.\hskip 1em(#1)}}\par}
\def\subproblem#1{\par\hangindent=1em\hangafter=0{%
\addtocounter{subprobcount}{1}\thesubproblem\emph{#1}\hskip 1em}}
\def\probskip{\vskip 10pt}
\def\medprobskip{\vskip 2in}
\def\subprobskip{\vskip 45pt}
\def\bigprobskip{\vskip 4in}

\begin{document}
{\emph{\fontsize{26}{28}\selectfont Math F252\hfill
{\fontsize{32}{36}\selectfont Midterm 2}
\hfill Fall 2025}}
\vskip 2cm
\strut\vtop{\halign{\emph#\hskip 0.5em\hfil&#\hbox to 2in{\hrulefill}\cr
\emph{\fontsize{18}{22}\selectfont Name:}&\cr
\noalign{\vskip 10pt}}}
%%\emph{\fontsize{18}{22}\selectfont Student Id:}&\cr
%%\noalign{\vskip 10pt}
%%\emph{\fontsize{18}{22}\selectfont Calculator Model:}&\cr
%}}
%\hfill
%\vtop{\halign{\emph{\fontsize{18}{22}\selectfont #}\hfil& \emph{\fontsize{18}{22}\selectfont\hskip 0.5ex $\square$ #}\hfil\cr
%Section: & 001 (Jill Faudree)\cr
%\noalign{\vskip 4pt}
%         & 002 (James Gossell)\cr
%\noalign{\vskip 4pt}
%         & 005 (James Gossell)\cr}}

\vfill
{\fontsize{18}{22}\selectfont\emph{Rules:}}

You have {90 minutes} to complete this midterm.  

Partial credit will be awarded, but you must show your work.

Calculators are not allowed. 

One hand-written sheet of notes is allowed.

%Place a box around your  \fbox{FINAL ANSWER} to each question where appropriate.

%If you need extra space, you can use the back sides of the pages.
%Please make it obvious  when you have done so.

Turn off anything that might go beep during the exam.

Good luck!
\vfill
\def\emptybox{\hbox to 2em{\vrule height 16pt depth 8pt width 0pt\hfil}}
\def\tline{\noalign{\hrule}}
\centerline{\vbox{\offinterlineskip
{
\bf\sf\fontsize{18pt}{22pt}\selectfont
\hrule
\halign{
\vrule#&\strut\quad\hfil#\hfil\quad&\vrule#&\quad\hfil#\hfil\quad
&\vrule#&\quad\hfil#\hfil\quad&\vrule#\cr
height 3pt&\omit&&\omit&&\omit&\cr
&Problem&&Possible&&Score&\cr\tline
height 3pt&\omit&&\omit&&\omit&\cr
&1&&8&&\emptybox&\cr\tline
&2&&8&&\emptybox&\cr\tline
&3&&12&&\emptybox&\cr\tline
&4&&8&&\emptybox&\cr\tline
&5&&24&&\emptybox&\cr\tline
&6&&8&&\emptybox&\cr\tline
&7&&12&&\emptybox&\cr\tline
&8&&12&&\emptybox&\cr\tline
&9&&8&&\emptybox&\cr\tline
&Extra Credit&&5&&\emptybox&\cr\tline
&Total&&100&&\emptybox&\cr
}\hrule}}}

\newpage
\begin{enumerate}
\item (8 pts) Consider the sequence $S=\left\{\frac{1}{2},\frac{2}{3},\frac{3}{4}, \frac{4}{5}, \frac{5}{6},\cdots \right\}.$
	\begin{enumerate}
	\item (4 pts) Find a formula for the general term $a_n$ of the sequence assuming that the pattern of the first few terms continues.
	\vfill
	\item (2 pts) Does this sequence converge? Justify your conclusion.
	\vfill
	\item (2 pts) Does series $\ds \sum_{n=1}^\infty a_n,$ with terms from the sequence $S$, converge? Justify your conclusion.
	\vfill
	\end{enumerate}
\item (8 pts) Determine if the integral below converges or diverges. Evaluate the integral if it converges. To earn full points, a solution must contain clear complete work and correct use of notation. \\
\quad \\
$\ds \int_0^4 \frac{1}{(4-x)^{2/3}} \: dx$

\vfill

\vfill


\newpage
\item (12 pts) For each \textbf{convergent} series below, determine its sum.  
	\begin{enumerate}
	\item $\ds \sum_{n=1}^{\infty} \left( \frac{3}{n^4}-\frac{3}{(n+1)^4}\right)$
	\vfill
	\item $\ds \sum_{n=1}^\infty \frac{(-5)^{n-1}}{2^{3n}}.$
	\vfill
	\end{enumerate}
\newpage
\item (8 pts) Use the \textbf{Integral Test} to determine if the series $\ds \sum_{n=1}^{\infty} ne^{-n^2}$ converges or diverges. (You do not have to verify that the Intergral Test applies.)

\vfill

\item (6 pts each) For each series below, show whether the series converges or diverges using an appropriate test. \textbf{State the test you use.}

	\begin{enumerate}
	\item $\ds \sum_{n=1}^{\infty} \frac{n^6}{2^n} $

	\vfill
	\newpage
	\item $\ds \sum_{n=1}^{\infty} \frac{n!}{(n+3)!} $

	\vfill

	\item $\ds \sum_{n=1}^{\infty} (-1)^{n+1}\ln\left(\tfrac{1}{n}\right)$ 

	\vfill

	
\newpage	
	\item $\ds \sum_{n=1}^{\infty} (-1)^{n+1} \frac{1}{\sqrt[3]{n+2}}$ 

	\vfill
	\end{enumerate}

\item (8 pts) Determine whether the series $\ds \sum_{n=1}^{\infty}  \frac{(-1)^{n+1}n^3}{5+n^5}$ is \textbf{absolutely convergent, conditionally convergent, or divergent.}
\vfill

\vfill




\newpage
\item (12 pts) Determine the radius of convergence, $R$, and the interval of convergence for each power series below.
	\begin{enumerate}
	\item  $\ds \sum_{n=0}^\infty \frac{(x-2)^n}{5n+4}$
	\vfill
	\vfill
	\item  $\ds \sum_{n=0}^\infty \frac{x^n}{n!}$
	\vfill
	\end{enumerate}
\newpage
\item (12 pts)
	\begin{enumerate}
	\item Find a power series representation for the function $\ds f(x)=\frac{x}{1-6x}.$
	\vfill
	\item Determine the radius of convergence and interval of convergence for the power series in part (a).
	\vfill
	\item Use your answer from part (a) to find a power series representation for $f'(x).$ 
	\vfill
	\end{enumerate}
\newpage
\item (8 pts) Find the Taylor series for $f(x)=\frac{1}{x}$ at $a=3.$ Your answer should be simplified.
	\vfill
	\vfill

\end{enumerate}

\textbf{Extra Credit} (5 pts) Determine the convergence of the two series below.
	\begin{enumerate}
	\item[a.]  $\ds \sum_{n=1}^\infty (\sqrt[n]{2} -1)^n$
	\vfill
	\item[b.]  $\ds \sum_{n=1}^\infty (\sqrt[n]{2} -1)$
	\vfill
	\end{enumerate}
\end{document}
